\chapter{Black Eye (Periorbital Ecchymosis)}

\section*{Definition}
Bruising around the eye caused by blunt trauma to the periorbital region, resulting in blood accumulation in the loose subcutaneous connective tissue of the eyelids.

\section*{When to See a Doctor}

\begin{itemize}
 \item Blurry or double vision, loss of vision, severe eye pain, pain or inability to move the eye, blood in the eye, an irregular pupil, or a high-velocity injury
 \item Bruising around both eyes after a head injury, loss of consciousness, vomiting, or a black eye without a direct injury
 \item Increasing warmth, pus, fever, or a black eye that has not improved within about three weeks
\end{itemize}

\section*{How It Develops}
Blunt force fractures small blood vessels in the periorbital tissues, causing blood to collect in the soft tissues. The loose eyelid skin allows significant swelling. Blood can track along tissue planes, sometimes spreading to the cheek or the other eye. The color changes as blood products break down, and uncomplicated bruising commonly improves over two to three weeks.

\section*{Causes}
\begin{itemize}
 \item Blunt trauma (sports injury, assault, fall)
 \item Nasal fracture with periorbital tracking
 \item Basilar skull fracture (raccoon eyes — bilateral, delayed)
 \item Orbital blowout fracture
 \item Post-surgical (eyelid surgery, rhinoplasty, sinus surgery)
 \item Coagulation disorders or anticoagulant medications
 \item Subgaleal hematoma tracking from scalp injury
 \item Insect bite or allergic reaction, which can mimic a black eye
 \item Periorbital cellulitis (not traumatic, requires urgent care)
 \item Child abuse (suspicious pattern or mechanism)
\end{itemize}

\section*{Related Symptoms}
\begin{itemize}
 \item Periorbital swelling
 \item Vision changes
 \item Headache
 \item Nasal pain
 \item Nausea
\end{itemize}

\section*{Medical Evaluation}
\begin{itemize}
 \item A clinician assesses the mechanism and timing of injury, visual acuity, pupils, eye movements, facial sensation, and signs of head injury.
 \item Slit-lamp examination is performed to evaluate for hyphema (blood in the anterior chamber), corneal abrasion, or traumatic iritis.
 \item CT of the orbits and/or head is considered when an orbital or facial fracture, foreign body, significant head injury, or other serious injury is suspected.
 \item Referral to an ophthalmologist is arranged for any associated visual disturbance, persistent diplopia, or suspicion of open-globe injury.
\end{itemize}

\section*{Treatment Options}
\begin{itemize}
 \item Most uncomplicated black eyes improve within two to three weeks with cold compresses initially followed by warm compresses.
 \item Orbital fractures are managed according to vision, eye movement, muscle entrapment, globe position, associated injuries, and symptoms; some need observation and specialist follow-up, while selected injuries need surgery.
 \item Corneal abrasions need an eye examination. A clinician may prescribe treatment and advise avoiding contact lenses until the surface has healed.
 \item People taking blood thinners or with a bleeding disorder should seek medical advice and should not stop prescribed medicine without direction; testing or treatment depends on the medication and findings.
\end{itemize}

\section*{Self-Care and Home Management}
\begin{itemize}
 \item Apply a cold compress (ice wrapped in a cloth) to the periorbital area for 15–20 minutes every 1–2 hours during the first 24–48 hours to reduce swelling.
 \item After 48 hours, switch to warm compresses to promote blood reabsorption and speed resolution of the bruise.
 \item Keep the head elevated on pillows while sleeping to reduce periorbital oedema.
 \item Avoid strenuous activity, bending over, or heavy lifting for the first few days as these may increase swelling and discoloration.
\end{itemize}

\section*{References}
\begin{thebibliography}{99}
\bibitem{black_eye:ref1} Bord SP, Linden J. ``Trauma to the globe and orbit.'' \textit{Emerg Med Clin North Am}. 2008;26(1):97-123.
\bibitem{black_eye:ref2} Jamal BT, Pfahler SM, et al. ``Blowout fractures: diagnosis and treatment.'' \textit{J Craniofac Surg}. 2009;20(1):182-185.
\bibitem{black_eye:ref3} American Academy of Ophthalmology. ``Preferred Practice Pattern: orbital trauma.'' AAO, 2021.
\bibitem{black_eye:ref4} Kanski JJ, Bowling B. \textit{Kanski\'s Clinical Ophthalmology}, 9th ed. Elsevier, 2020.
\bibitem{black_eye:ref5} National Health Service. ``Black eye.'' Reviewed July 2023. https://www.nhs.uk/conditions/black-eye/.
\bibitem{black_eye:ref6} American Academy of Ophthalmology. ``Orbital Floor Fracture.'' EyeWiki, accessed September 2026. https://eyewiki.aao.org/Orbital\_Floor\_Fracture.
\end{thebibliography}
